\chapter{Conclusions and future work}\label{chap:conclusions}
\thispagestyle{fancy} %Resets custom page number location on first chapter page

This thesis asked whether the positions of a telescope's mirrors can be inferred directly from images of point sources. We presented a Bayesian pipeline that does this for simulated multi-field observations of telescopes with known prescriptions. A ray-trace model in Optiland applies rigid-body misalignments to the mirrors and computes the PSF at several field points. Each PSF is reduced to a short vector of basis coefficients, and a multivariate Gaussian likelihood compares these coefficients with the observed ones. Nested sampling with Nautilus explores the posterior over the alignment parameters, and differential evolution, started from and bounded by that posterior, refines a point estimate. This chapter summarizes the results (Section~\ref{sec:conc_results}), discusses what they show about the two stages of the pipeline (Section~\ref{sec:conc_stages}) and about the limits of its precision (Section~\ref{sec:conc_precision}), states the limitations of the study (Section~\ref{sec:conc_limits}) and outlines future work (Section~\ref{sec:conc_future}).

\section{Summary of results}\label{sec:conc_results}

Two benchmarks tested the pipeline (Table~\ref{tab:conc_summary}). In both, the alignment was inferred from noisy observations, with Gaussian noise of $10^{-4}$ of the ideal PSF peak added to every observed image.

The first benchmark was a Ritchey--Chr{\'e}tien model of the Hubble Space Telescope with five secondary-mirror degrees of freedom, observed at five field points and described by 15 circular-Bessel coefficients per field (Chapter~\ref{chap:hubble}). Nested sampling constrained the despace and both tilts but left the two decenters nearly unconstrained across the $\pm200\,\mu\mathrm{m}$ prior. Differential evolution, searching within bounds taken from the posterior, recovered the translations to within $4\,\mu\mathrm{m}$ and the tilts to within about $4\times10^{-4}$ degrees. It reduced the on-axis RMS wavefront error from 2.26 waves to $1.3\times10^{-3}$ waves.

The second benchmark was a three-mirror anastigmat with ten degrees of freedom on its primary and secondary mirrors, observed at ten field points and described by 11 Zernike coefficients per field (Chapter~\ref{chap:tma}). Its posterior was broader. The primary decenter $pdx$ and both primary tilts were nearly unconstrained, the secondary decenters were broad, and several pairs of parameters were strongly correlated. Differential evolution recovered five of the six translations to within $10.3\,\mu\mathrm{m}$ and the secondary decenter $sdx$ to within $46\,\mu\mathrm{m}$. It reduced the on-axis RMS wavefront error from 8.24 waves to $2.5\times10^{-3}$ waves.

Chapter~\ref{chap:noise} showed that these residuals are the ones the noise predicts. Both differential-evolution runs reached the noise floor of their objective, and every parameter residual lies within 1.6 predicted standard deviations of zero. The final wavefront errors also lie within the range obtained by propagating the noise through the linearized forward model. The remaining errors therefore come from the noise in the observations, not from the optimizer.

Chapter~\ref{chap:noise} also tested a way to improve on this limit for Hubble: observing the secondary again after known moves (Section~\ref{sec:noise_moves}). With four tilts of $\pm0.1$ degrees, the predicted noise-limited errors were 6--7 times smaller than from the nominal position alone, and 16 times smaller on the despace; a factor of $\sqrt{5}$ comes from the four additional exposures and the rest from the moves. A full run of the pipeline recovered both decenters to within $0.2\,\mu\mathrm{m}$ and reduced the on-axis RMS wavefront error to $2.1\times10^{-4}$ waves.

\begin{table}
\centering
\caption[Summary of the two benchmarks.]{Summary of the two benchmarks. WFE is the on-axis RMS wavefront error of the system with residual misalignment $\boldsymbol{\theta}_{\mathrm{true}}-\mathbf{c}$ (Section~\ref{sec:raytracing}). The noise-limited range is the 16th--84th percentile range predicted from the noise alone (Section~\ref{sec:noise_de}).}
\label{tab:conc_summary}
\begin{tabular}{lll}
\hline
 & \textbf{Hubble} & \textbf{TMA} \\
\hline
Mirrors perturbed & secondary & primary, secondary \\
Degrees of freedom & 5 & 10 \\
Fields $\times$ modes & $5\times15$ & $10\times11$ \\
Likelihood evaluations (nested sampling) & 44\,000 & 93\,200 \\
Evaluations of $S$ (DE, before polishing) & 150\,075 & 300\,150 \\
WFE, uncorrected (waves) & 2.26 & 8.24 \\
WFE, nested-sampling start $\mathbf{x}_0$ (waves) & 0.014 & 0.99 \\
WFE, DE solution (waves) & $1.3\times10^{-3}$ & $2.5\times10^{-3}$ \\
WFE, noise-limited range (waves) & $(0.7$--$1.7)\times10^{-3}$ & $(2.4$--$6.9)\times10^{-3}$ \\
Largest translation residual ($\mu$m) & 3.8 ($sdy$) & 46 ($sdx$) \\
Largest tilt residual (degree) & $4.1\times10^{-4}$ ($stx$) & $3.7\times10^{-4}$ ($pty$) \\
\hline
\end{tabular}
\end{table}

\section{The roles of the two stages}\label{sec:conc_stages}

The two stages of the pipeline do different jobs, and the benchmarks show that both are needed. Nested sampling is the global search. Starting from the full prior, it found the high-likelihood region in both benchmarks, and its starting point alone reduced the on-axis wavefront error by a factor of about 160 for Hubble and about 8 for the TMA. The posterior also maps the structure of the problem. Its marginals show which parameters the observations constrain, and its correlations reveal combinations of mirror motions that leave the PSFs nearly unchanged. An optimizer alone would return a single point without this information. Nested sampling also computes the evidence, which we report (Table~\ref{tab:nautilus}) but have not yet used.

The posterior does not, however, give the precision of the measurement. Its width is set by the calibration covariances $\boldsymbol{\Sigma}_f$, which describe how the coefficients vary across the prior rather than the measurement noise; at the noise level of the benchmarks, the noise shifts $-2\ln\mathcal{L}$ by less than 0.03 at the truth (Section~\ref{sec:noise_likelihood}). Differential evolution supplies the precision. By minimizing the unweighted coefficient residuals, it reduced the wavefront error by a further factor of about 10 for Hubble and about 400 for the TMA and reached the noise floor in both runs. Over both stages, the wavefront error fell by factors of about 1700 and 3300.

Nested sampling also made this refinement tractable. The bounds taken from the innermost posterior shell narrowed the search in the well-constrained parameters by more than an order of magnitude, for example from $\pm200\,\mu\mathrm{m}$ to 36--$43\,\mu\mathrm{m}$ on the Hubble despace and from $\pm0.4$ degrees to ranges 0.02--0.03 degrees wide on the TMA secondary tilts, and contained the true alignment in both runs. L-BFGS-B polishing left both solutions unchanged, but this says little about convergence, because its finite-difference gradients are dominated by numerical noise in the objective (Section~\ref{sec:de}). That both objectives ended at their predicted noise floors (Section~\ref{sec:noise_de}) is what shows that differential evolution reached the minimum.

\section{What limits the precision}\label{sec:conc_precision}

The precision of the recovered alignment is set by the noise and by the near-degeneracies of the problem together. Near-degenerate combinations of mirror motions appeared in both benchmarks, and they have a simple optical origin. A decenter of a mirror and a tilt about the perpendicular axis, of the same mirror or of another, both introduce coma that is nearly constant across the field, so one can largely cancel the other. This links $sdy$ with $stx$ and $sdx$ with $sty$ for Hubble, and $sdy$ with $ptx$ and $sdx$ with $pty$ and $sty$ for the TMA. The two despaces of the TMA trade off against each other in setting focus. These compensating motions are the ones described by nodal aberration theory~\citep{Thompson2005,SchechterLevinson2011}. A compensating motion cancels the field-constant aberration but leaves a field-dependent residual, which is why the likelihood must combine several field points.

The linearized analysis of Chapter~\ref{chap:noise} quantified what these degeneracies cost. If all other parameters were known, the noise alone would determine the Hubble decenters to $0.36$ and $0.32\,\mu\mathrm{m}$ and the TMA secondary decenter $sdx$ to $0.30\,\mu\mathrm{m}$. The correlations amplify these errors by a factor of about 10 for Hubble and about 170 for the TMA. The largest residual of the TMA run, $46\,\mu\mathrm{m}$ on $sdx$, is therefore noise projected onto a nearly degenerate direction rather than a failure of the optimizer.

Because the degenerate directions change the PSFs little, errors along them also change the image quality little: the TMA solution, with its $46\,\mu\mathrm{m}$ error in $sdx$, leaves only $2.5\times10^{-3}$ waves of on-axis wavefront error. The two Hubble runs with known moves show the same effect: their decenter errors differ by more than a factor of ten, yet they leave the same on-axis wavefront error, about $2\times10^{-4}$ waves (Section~\ref{sec:noise_moves}). For correcting the image quality at the observed fields, the degeneracies matter less than the parameter residuals suggest. They matter more when the mirror positions themselves are needed, for example to track mechanical drifts, or when the aberrations at field points that were not observed must also be controlled.

The noise sets the scale of these errors. At the noise level of the benchmarks, the parameter errors and the wavefront error grow in proportion to the noise up to $3\times10^{-3}$ of the ideal peak for Hubble and up to $3\times10^{-4}$ for the TMA (Section~\ref{sec:noise_sweep}). Beyond that, the TMA despace errors grow faster than linearly as the trade-off between the two despaces becomes non-linear, and the linear prediction places the median wavefront error at the Mar{\'e}chal criterion of $\lambda/14$ near $1.7\times10^{-3}$ of the peak. Within the assumptions of this study, the pipeline could therefore bring a TMA of this design to a median wavefront error within the Mar{\'e}chal criterion from observations whose noise is below about $0.17\%$ of the ideal PSF peak.

The precision can also be raised by making the observations more informative rather than less noisy. Moving the Hubble secondary by known tilts and observing again adds known coma to the PSFs, which makes the coefficients respond more strongly to every parameter, much as the known defocus of phase diversity does (Section~\ref{sec:noise_moves}). With tilts of $\pm0.1$ degrees, the predicted errors fell by a further factor of about 3 on the decenters and tilts, and 7 on the despace, beyond the $\sqrt{5}$ gained from the additional exposures. The moves did not, however, weaken the degeneracies: the decenter--tilt correlations stayed above 0.98 in magnitude. Their benefit depends on their size, since a move must be comparable to the coma already present to change the PSF appreciably; tilts of $\pm0.05$ degrees gave a further factor of only 1.3--1.4, and tilts of $0.01$ degrees or less none. It also depends on knowing the moves, to about $0.8\,\mu\mathrm{m}$ in decenter, $5\times10^{-5}$ degrees in tilt and $5$\,nm in despace per move for the $\pm0.1$-degree design.

\section{Limitations}\label{sec:conc_limits}

The benchmarks are idealized in several ways, and the results should be read with these in mind.
\begin{itemize}
    \item \textbf{The forward model is exact.} The observations were generated with the same Optiland model used for inference, so the prescriptions, the aperture masks and the positions of the fixed optics and the detector are known exactly. A real telescope differs from its design through figure errors on the mirrors, uncertainties in their radii and conic constants, and misplaced fixed components. Any such mismatch would bias the inferred positions, and its effect has not been measured here.
    \item \textbf{The PSFs are simplified.} They are monochromatic, sampled at 4 and 8 pixels per $\lambda F_\#$, and free of detector effects such as pixel integration and jitter. They are also centered on the chief ray, so the position of the image on the detector is not used. The noise is additive, Gaussian and independent of the signal; photon noise, background and detector systematics are not modeled.
    \item \textbf{Each benchmark is a single case.} Each telescope was tested with one injected misalignment and one noise realization per field. Chapter~\ref{chap:noise} shows that the residuals are typical of the noise, but the success rate of the pipeline over many misalignments, in particular large ones near the edges of the prior, has not been measured.
    \item \textbf{The posterior widths are not measurement uncertainties.} Because the likelihood uses the calibration covariances $\boldsymbol{\Sigma}_f$, the credible intervals in Figures~\ref{fig:hubble_corner} and~\ref{fig:pearl_corner} describe how the coefficients vary across the prior; the uncertainty of the final estimate comes from the linearized analysis of Chapter~\ref{chap:noise}, not from the sampler. The Hubble prior transform also rounds the tilts to 0.01 degrees, far coarser than the precision that differential evolution reached. In the runs with known moves of the secondary, this rounding split the posterior into separate modes about $100\,\mu\mathrm{m}$ apart in the decenters and, with tilts of $\pm0.1$ degrees, collapsed the differential-evolution bounds on the tilts onto single values, which had to be widened by hand (Section~\ref{sec:noise_moves}).
    \item \textbf{Only some mirrors move.} The Hubble model perturbs only the secondary, and the TMA only the primary and secondary. The tertiary, the fold mirror and the detector stay fixed, and rotations about the optical axis are not included.
    \item \textbf{The computation is expensive.} Counting one PSF per field for every likelihood evaluation and every evaluation of the DE objective, the Hubble benchmark computed about $10^6$ PSFs and the TMA about $4\times10^6$, roughly three quarters of them during differential evolution.
\end{itemize}

\section{Future work}\label{sec:conc_future}

\subsection{Realistic observations}
The next step toward applying the pipeline is to make the simulated observations realistic and then to replace them with measured PSFs. The forward model can be extended to broadband PSFs by summing monochromatic PSFs over the passband, and to real detectors by integrating over pixels at the plate scale of an instrument. The noise model should include photon noise, read noise and background. For bright stars, photon noise is close to Gaussian with a variance that changes across the image; because the basis projection is linear, the resulting coefficient covariance follows from equation~(\ref{eq:noisecov}) with $\sigma^2$ moved inside the sum over pixels. The effect of model mismatch can be measured by injecting errors that the inference model does not contain, such as mirror figure errors or small changes in the conic constants, and recording the bias they cause in the inferred positions. The most important of these errors could then be added to the model as nuisance parameters and marginalized over. Finally, the pipeline should be applied to laboratory testbeds with known, actuated misalignments and then to on-sky PSF campaigns relevant to Habitable Worlds Observatory class TMA designs~\citep{HWO2023}.

\subsection{A likelihood based on the noise}
At present, the posterior locates the solution and differential evolution measures it. A likelihood built from the noise covariance $\mathbf{N}$ would instead give posterior widths that are measurement uncertainties, with the correlations of Figure~\ref{fig:noise_ellipses} built in. Such a likelihood is far more sharply peaked, so a practical scheme would keep the present nested-sampling stage as the global search and sample the noise-based likelihood locally, for example with Markov chain Monte Carlo started from the differential-evolution solution. At the benchmark noise level, the linear predictions of Chapter~\ref{chap:noise} held to within about 20\%, which suggests that this posterior would be close to the Gaussian of equation~(\ref{eq:lscov}). A likelihood that describes the measurement would also make the evidence useful: comparing models in which different subsets of the degrees of freedom are free could identify which mirrors moved.

\subsection{Breaking the degeneracies}
The linearized analysis of Chapter~\ref{chap:noise} can also guide the design of observations. The covariance $\mathbf{A}\,\mathbf{N}_{\mathrm{all}}\,\mathbf{A}^{\top}$ requires only the Jacobian of the coefficients at a nominal alignment, so it can be evaluated for any candidate set of field points, basis or noise level before data are taken. It could be used to choose field points that best separate a decenter from a compensating tilt; the field-dependent aberrations that distinguish them grow with field angle. Other observables could break the degeneracies more directly: the position of each image on the detector, which the present pipeline discards because its PSFs are centered on the chief ray, although it moves by millimeters along the Hubble decenter--tilt degeneracy (Section~\ref{sec:noise_moves}); defocused or phase-diverse images~\citep{Gonsalves1982}; PSFs at several wavelengths; or more basis modes, which would retain information that the present 11- and 15-mode projections discard. Section~\ref{sec:noise_moves} tested one form of diversity, PSFs recorded after known moves of the secondary, and chose the size of the moves with this calculation. The moves raised the sensitivity of the coefficients severalfold but left the decenter--tilt correlations unchanged, so diversity of this kind complements, rather than replaces, observables that separate a decenter from a compensating tilt. For the TMA, where the degeneracies amplify the noise by more than 100, the same calculation could identify which moves of which mirror would help most. Where they are available, actuator or metrology readings could also constrain individual mirrors through the prior. Each of these options can be assessed with the same linear calculation before it is added to the full pipeline.

\subsection{More degrees of freedom}
Aligning a real observatory involves more degrees of freedom than the benchmarks: the tertiary and fold mirrors, the detector position, rotations about the optical axis for optics that are not rotationally symmetric, low-order changes in the figure of each mirror and, for segmented primaries such as those of JWST and several HWO concepts, the piston and tip-tilt of every segment. Doubling the number of parameters from Hubble to the TMA already broadened the posterior, strengthened the degeneracies and roughly doubled the number of likelihood evaluations. Scaling to tens or hundreds of parameters will need the efficiency gains described below and the degeneracy breaking described above. A staged approach, in which the largest motions are solved first and finer terms are added around that solution, would follow the sequence of JWST commissioning~\citep{Acton2022}.

\subsection{Computational cost}
Several changes could reduce the cost of the pipeline. Neither differential-evolution run met its convergence tolerance, because the population approaches the noise floor only slowly, and both ran the full 2000 generations although the Hubble objective was within 4\% of its final value by generation 30. In the linear approximation, a parameter error of one standard deviation, in any direction, raises the objective by about the noise variance of a single coefficient. A stopping rule that compares recent improvements of the objective with this variance would end a run once further generations could no longer change the estimate meaningfully.

Differential evolution could also be replaced, after nested sampling, by a local least-squares solver. The Monte Carlo study of Section~\ref{sec:noise_sweep} solved the same least-squares problem with twelve Gauss--Newton iterations started from the truth. In practice, each iteration would need a new Jacobian, or $2d+1$ forward-model evaluations by central differences, so a few tens of iterations cost hundreds of evaluations rather than the $1.5$--$3\times10^5$ used by differential evolution. A damped Gauss--Newton solver, as in merit-function regression~\citep{Kim2007}, started from the nested-sampling point could therefore reach the noise floor at a small fraction of the cost. A test on the Hubble benchmark, with a new noise realization, supports this. Started from the nested-sampling point, L-BFGS-B with gradients from Optiland's PyTorch backend came within 0.1 noise standard deviations of the differential-evolution solution after 31 evaluations of $S$, against about 5900 for differential evolution. Seven further starts spread over the nested-sampling bounds converged to the same minimum, and so did L-BFGS-B with central-difference gradients, with about ten times as many evaluations. SciPy's default forward-difference gradients failed from every start, for the reason given in Section~\ref{sec:de}. All the starts lay within the nested-sampling bounds, so nested sampling is still needed to provide them.

The forward model itself could be made faster. Every likelihood evaluation is independent, so both stages parallelize; the TMA run already evaluated each DE generation in parallel. Ray tracers with GPU and automatic-differentiation backends, including Optiland's PyTorch backend, could trace many misalignments at once and provide exact gradients for local refinement and for gradient-based samplers such as Hamiltonian Monte Carlo. In Optiland 0.6.0, however, the gradients with respect to decenters and tilts are correct only after a small fix. The reference sphere for the wavefront is built from detached scalar values, so automatic differentiation misses its dependence on the misalignment and returns gradients of the wrong sign that are 130--1900 times too large. Keeping the sphere in the computation graph makes them agree with central differences to three or more digits. Finally, the calibration step already computes $10^4$ forward-model evaluations per field. A neural-network emulator of the coefficients trained on such draws could replace the ray trace inside the sampler, with ray tracing kept for the final refinement and validation. Simulation-based inference could go further and learn the posterior as a function of the observed coefficients, so that each new observation could be analyzed without a new sampling run.

\subsection{Toward closed-loop alignment}
In operation, alignment is a sequence of measurements rather than a single one: a telescope drifts, for example as its temperature changes, and is corrected repeatedly. Each update could use the previous solution and its uncertainty as the prior, which would narrow the search and reduce its cost. Combining the inferred positions with a model of the mirror actuators would close the loop: observe, infer the mirror positions, command corrective moves and observe again. A laboratory testbed with actuated mirrors is the natural place to demonstrate this loop.

\section{Closing remarks}\label{sec:conc_closing}

Phase retrieval and wavefront sensing show that a telescope is misaligned; this work shows that, for simulated telescopes with known prescriptions, the misalignment itself can be inferred from multi-field PSFs. In both benchmarks the recovered alignment left a few thousandths of a wave of on-axis wavefront error, and the remaining errors in the mirror positions are those set by the noise and amplified along the near-degenerate directions of the problem. The approach is computationally intensive, but it measures the state of the optical system directly, in principle from a single image of a wide-field instrument, as is typical of modern TMA observatories, rather than through a linearized model of how the wavefront responds to mirror motions. With realistic observations, faster forward models, known mirror moves that sharpen the measurement and observables chosen to break the degeneracies, it could enable deterministic alignment of future observatories such as HWO~\citep{HWO2023} and Lazuli~\citep{Roy2026}.
